\clearpage
\VerdierPhysicalPageAnchor{0232}%
\begingroup\small\setlength{\parskip}{0pt}\setlength{\emergencystretch}{1.5em}
\setlength{\abovedisplayskip}{4pt plus 2pt minus 2pt}%
\setlength{\belowdisplayskip}{4pt plus 2pt minus 2pt}%

\VerdierSemanticAnchor{ch4d-4-1-7-double-complex-case}
Suppose then that $X$ is a double complex and that $m=1$. To prove a), it
suffices to show that, for every $(p,q,r)$ in
$\mathcal{F}\ell_2(\widetilde{\mathbb Z})$, the following diagram in
$\operatorname{comp}(\mathcal A)$:

\VerdierSemanticAnchor{ch4d-4-1-7-6}
\begin{equation}
\resizebox{0.98\linewidth}{!}{$
\displaystyle
\int_2{}^{1}\!X(p,q)\longrightarrow
\int_2{}^{1}\!X(p,r)\longrightarrow
\int_2{}^{1}\!X(q,r)
\xrightarrow{{}^{1}\!d(p,q,r)}
\int_2T_1\bigl({}^{1}\!X(p,q)\bigr)
=\left(\int_2{}^{1}\!X(p,q)\right)[1]
$}
\tag{4.1.7.6}
\end{equation}

defines a distinguished triangle in $K(\mathcal A)$ (observe that in this
case the commutation isomorphism for translations
$\int_2T_1({}^{1}\!X(p,q))\xrightarrow{\sim}
(\int_2{}^{1}\!X(p,q))[1]$ is the identity). The other assertions in a)
are evident. To prove that diagram (4.1.7.6) defines a distinguished
triangle in $K(\mathcal A)$, it suffices, by property (TRII) (Chapter~II,
1.1.1), to show that the diagram

\VerdierSemanticAnchor{ch4d-4-1-7-7}
\begin{equation}
\resizebox{0.96\linewidth}{!}{$
\displaystyle
\left(\int_2{}^{1}\!X(q,r)\right)[-1]
\xrightarrow{-{}^{1}\!d(p,q,r)[-1]}
\int_2{}^{1}\!X(p,q)\longrightarrow
\int_2{}^{1}\!X(p,r)\longrightarrow
\int_2{}^{1}\!X(q,r)
$}
\tag{4.1.7.7}
\end{equation}

defines a distinguished triangle in $K(\mathcal A)$. One sees immediately
that the simple complex $\int_2{}^{1}\!X(p,r)$ is equal to the cone of the
morphism $-{}^{1}\!d(p,q,r)[-1]$ (Chapter~I, 3.1.2), and that the
morphisms
$\int_2{}^{1}\!X(p,q)\longrightarrow\int_2{}^{1}\!X(p,r)$ and
$\int_2{}^{1}\!X(p,r)\longrightarrow\int_2{}^{1}\!X(q,r)$ appearing in
(4.1.7.7) are the morphisms introduced in the study of the cone functor
(Chapter~I, 3.2.2.2 and 3.2.2.4). Consequently, the triangle in
$K(\mathcal A)$ defined by diagram (4.1.7.7) is, by definition
(Chapter~I, 3.3.1), a distinguished triangle. We now prove b). Using the
notation introduced in (4.1.6.5), we have just shown that in the diagram

\VerdierSemanticAnchor{ch4d-4-1-7-8}
\begin{equation}
{}^{1}\!Y(q,r)[-1]
\xrightarrow{-{}^{1}\!d(-\infty,q,r)}
{}^{1}\!Y(q)\longrightarrow{}^{1}\!Y(r)\longrightarrow{}^{1}\!Y(q,r),
\tag{4.1.7.8}
\end{equation}

the complex ${}^{1}\!Y(r)$ appears as the cone of the morphism
$-{}^{1}\!d(-\infty,q,r)$ (set $p=-\infty$ in diagram (4.1.7.7)).
Moreover, in (4.1.1), the object $\mathcal Y(q,r)$ was defined as the cone
of the morphism

\VerdierSemanticAnchor{ch4d-4-1-7-9}
\begin{equation}
{}^{1}\!Y(q)\longrightarrow{}^{1}\!Y(r).
\tag{4.1.7.9}
\end{equation}

In this situation (Chapter~I, 3.2.5.2), a homotopy equivalence (and hence
an isomorphism in $K(\mathcal A)$) was then defined:

\VerdierSemanticAnchor{ch4d-4-1-7-10}
\begin{equation}
J(q,r)\colon\mathcal Y(q,r)\longrightarrow{}^{1}\!Y(q,r).
\tag{4.1.7.10}
\end{equation}

\VerdierSemanticAnchor{ch4d-4-1-7-proof-end}
We leave it to the reader to verify that this homotopy equivalence is
functorial with respect to the indices $(q,r)$ in
$\mathcal{F}\ell_1(\widetilde{\mathbb Z})$ and with respect to the double
complex $X$, and that it defines a canonical isomorphism from the spectral
object $\bigl(\mathcal Y(p,q),\delta(p,q,r)\bigr)$ onto the spectral object
$\bigl({}^{1}\!Y(p,q),{}^{1}\!d(p,q,r)\bigr)$.
\endgroup
